

Or si \(K\) est un \(\lambda\)-anneau, le diagramme :
\begin{equation}\tag{7.5.4}
\begin{tikzcd}[column sep=large,row sep=large]
K \arrow[r,"\lambda_t"] \arrow[d,"\psi^{k'}"'] &
\widehat G(K) \arrow[d,"\psi^{k'}"]\\
K \arrow[r,"\lambda_t"] &
\widehat G(K)
\end{tikzcd}
\end{equation}
est commutatif d'après 7.3, puisque \(\lambda_t\) est un \(\lambda\)-homomorphisme. On en déduit aussitôt les relations cherchées.

ii) Si \(K\) est de caractéristique nulle, l'application « dérivée logarithmique » est un isomorphisme de \(\widehat G(K)\) sur \(K[[t]]\), car on peut intégrer. Par suite, deux éléments de \(\widehat G(K)\) sont égaux si, et seulement si pour tout \(k\), ils ont même image par \(\Phi^k\).

Pour montrer que \(\lambda_t\) est un homomorphisme d'anneaux, il faut montrer la commutativité du carré de gauche du diagramme (7.5.2), et d'après la remarque précédente, il suffit de montrer que pour tout \(k\), le diagramme (7.5.2) est commutatif, puisque le carré de droite l'est toujours. Or ceci n'est autre que les relations \(\psi^k(xy)=\psi^k(x)\psi^k(y)\).

Pour montrer que \(\lambda_t\) est un \(\lambda\)-homomorphisme, il suffit de montrer que le diagramme (7.5.4) est commutatif pour tout \(k'\), d'après 7.3. Or, d'après la remarque qui précède, cela équivaut à montrer pour tous \(k,k'\), que \(\Phi^k\circ\psi^{k'}\circ\lambda_t=\Phi^k\circ\lambda_t\circ\psi^{k'}\), c'est-à-dire les relations \(\psi^{kk'}=\psi^k\circ\psi^{k'}\).

Enfin, la relation \(\lambda_t(1)=1+t\) s'obtient par intégration à partir des relations \(\psi^k(1)=1\).

\subsection*{BIBLIOGRAPHIE}
\begin{enumerate}[label={[\arabic*]}]
\item A. Grothendieck, La théorie des classes de Chern, Bull. Soc. math. France, 86, 1958, p. 137 à 154.
\item F. Hirzebruch, Neue topologische Methoden in der algebraischen Geometrie.
\end{enumerate}



% Exposé VI projective-bundle macro additions
\providecommand{\Sym}{\operatorname{Sym}}
\providecommand{\Proj}{\operatorname{Proj}}
\providecommand{\Projz}{\operatorname{Proj}_0}
\providecommand{\Spec}{\operatorname{Spec}}
\providecommand{\Comp}{\operatorname{Comp}}
\providecommand{\Dplus}{D_{+}}
\renewcommand{\id}{\operatorname{id}}
\providecommand{\cG}{\mathcal G}
\providecommand{\cA}{\mathcal A}
\providecommand{\cB}{\mathcal B}
\providecommand{\cL}{\mathscr L}
\providecommand{\cM}{\mathscr M}
\providecommand{\cN}{\mathscr N}
\providecommand{\cE}{\mathscr E}
\providecommand{\cK}{\mathscr K}
\providecommand{\cP}{\mathscr P}
\providecommand{\cQ}{\mathscr Q}
\providecommand{\OO}{\mathcal O}
\providecommand{\rk}{\operatorname{rk}}
\providecommand{\colim}{\operatorname*{colim}}



